% ============================================================
% Section 1: Introduction
% ============================================================

\section{Introduction}
\label{sec:intro}

Optical speckle is the high-contrast granular pattern that emerges when coherent
light, such as that of a laser, is scattered by a rough surface.
This phenomenon appears in optical metrology, optical coherence tomography,
biomedical imaging, optical cryptography and fiber communications
\cite{Goodman2007_Speckle}.
Its accurate prediction requires solving the Helmholtz equation with stochastic
boundary conditions, a problem that classical numerical methods address at a
prohibitive computational cost.

Mesh-based methods, such as the Finite Element Method (FEM), are accurate but
impose severe constraints: the mesh must have a spatial resolution finer than the
laser wavelength in order to capture the fine details of the speckle pattern
\cite{Jin2014_FEM}.
This leads to systems with millions of unknowns per realization, and every new
surface roughness profile requires solving the complete system from scratch.
For statistical studies requiring thousands of realizations, the cost becomes
intractable.

Physics-Informed Neural Networks (PINNs), introduced by Raissi et al.
\cite{Raissi2019_PINNs}, propose an alternative \textit{mesh-free} paradigm:
embedding the physical laws directly into the loss function of the network, so
that the fit to the boundary conditions and the fulfillment of the differential
equation over the domain are enforced simultaneously.
Once trained, the model evaluates the complete field on a $100 \times 100$ grid in
$\approx 1.3\,\mathrm{ms}$ (NVIDIA RTX~5050 GPU), acting as a computationally
efficient surrogate of the numerical solver.
The effectiveness of this paradigm has been demonstrated in wave scattering
problems \cite{Fang2020_scattering}, nano-optics and parameter inversion
\cite{Chen2020_nanoptics}, and the PINN framework has been consolidated in the
reviews by Karniadakis et al. \cite{Karniadakis2021_review} and Cuomo et al.
\cite{Cuomo2022_review}.

A critical obstacle of standard PINNs for the Helmholtz equation is the
instability of second-order derivatives under ReLU or tanh activations.
Sinusoidal Representation Networks (SIREN), proposed by Sitzmann et al.
\cite{Sitzmann2020_SIREN}, address this problem through activations
$\phi(z) = \sin(\omega_0 z)$ together with an initialization scheme that preserves
the distribution of activations throughout the network, making the architecture a
natural choice for wave equations.
A recent study by Schoder and Kraxberger \cite{Schoder2024_PINN_acoustics}
validated the feasibility of PINNs for the 3D Helmholtz equation, reporting an
$L^2$ error of $2.5\%$ with respect to reference FEM solutions.
This paradigm has recently been applied to optical and electromagnetic settings
with a SIREN architecture \cite{Panagiotakopoulos2026_Helmholtz}.
However, to the best of our knowledge, no previous work combines sinusoidal
activation, an explicit calibration rule for $\omega_0$ with respect to the
wavenumber, and the generation of optical speckle from a random-phase boundary
condition statistically validated against Goodman's criterion.

In this work we present a PINN-SIREN implementation for solving the
two-dimensional Helmholtz equation with a complex field, laying the groundwork for
the future accelerated simulation of optical speckle.
The main contribution is the validation of the SIREN $5\times128$ architecture in
two scenarios of increasing complexity: 1D validation against an exact analytical
solution ($L^2$ error $= 0.006\%$) and 2D validation with a complex plane wave
($L^2$ error $= 0.171\%$), both several orders of magnitude below the $5\%$
threshold adopted in this work.
The generation of speckle patterns with a random-phase boundary condition
$\phi(x) \sim \mathcal{U}(0, 2\pi)$ constitutes the next planned stage of this
research.

The article is organized as follows.
Section~\ref{sec:theory} presents the theoretical framework: the Helmholtz
equation, the PINN formulation, the SIREN architecture and speckle statistics.
Section~\ref{sec:methodology} describes the problem formulation, LHS sampling and
the Adam + L-BFGS optimization strategy.
Section~\ref{sec:results} reports the results of the 1D and 2D validation, together
with the comparison against the state of the art.
Section~\ref{sec:conclusions} summarizes the contributions and outlines future
directions.
